%% ============================================================
%% SECTION 9: FUTURE RESEARCH DIRECTIONS
%% EarthVision-LM Survey -- Artificial Intelligence Review
%% Template: sn-jnl (Springer Nature) | sn-mathphys-num
%% Self-citations framed as "promising direction" -- not "we prove"
%% Rules: booktabs, TikZ, numbered [N] cites, NO fabricated data
%% ============================================================

\section{Future Research Directions}\label{sec9}

The analyses of Sects.~\ref{sec3}--\ref{sec8} collectively expose a set of
structural gaps between current RS-MLLM capability and the requirements of
operational geospatial intelligence. This section first quantifies those
gaps through an Operational Readiness Matrix, then organises eight concrete
research directions grounded in the specific deficits identified throughout
the survey.

\subsection{Operational Readiness Matrix}\label{sec9:orm}

Table~\ref{tab:orm} assesses the current RS-MLLM field against operational
requirements across nine capability dimensions. The matrix makes the
``Road to Geospatial Intelligence'' argument of the title concrete and
falsifiable: each row identifies a measurable gap between where the field
stands today and what operational deployment requires.

\begin{table*}[t]
\caption{Operational Readiness Matrix for RS-MLLMs. \textbf{Current RS-MLLM}:
capability level across published models as of August 2026.
\textbf{Operational requirement}: capability level needed for reliable
geospatial deployment (disaster response, infrastructure monitoring,
agricultural management, border surveillance).
\textbf{Gap}: difference between current and required.
Ratings: Very High\,/\,High\,/\,Moderate\,/\,Low\,/\,Emerging.
Evidence tier per Sect.~\ref{sec2:prisma}: A\,=\,peer-reviewed controlled;
B\,=\,peer-reviewed benchmark; C\,=\,preprint; D\,=\,conceptual.
}\label{tab:orm}
\small\setlength\tabcolsep{3pt}
\begin{tabular*}{\textwidth}{@{\extracolsep\fill}p{2.4cm}p{2.6cm}p{2.8cm}p{1.6cm}p{3.2cm}p{0.7cm}}
\toprule
\textbf{Capability} & \textbf{Current RS-MLLM} & \textbf{Operational requirement} & \textbf{Gap} & \textbf{Key evidence} & \textbf{Tier} \\
\midrule
VQA (optical)        & Strong (95.2\,\% RSVQA-LR)    & Strong                  & Low         & Earth-OneVision RSVQA-LR score~\cite{earthonevision2026}             & C \\
Captioning           & Strong (CIDEr 91.4)            & Moderate                & Low         & Earth-OneVision RSICD CIDEr~\cite{earthonevision2026}                & C \\
HBB Grounding        & Moderate (Acc@0.5 84.6)        & High                    & Medium      & RSVG benchmark~\cite{rsvg2023}                                       & B \\
OBB Detection        & Weak (28.6 zero-shot mAP)      & Critical                & Very High   & Specialist: 75.9; RS-MLLM: 28.6~\cite{xie2021orientedrcnn,geochat2024} & A \\
SAR reasoning        & Weak (no physics encoder)      & Critical                & Very High   & No SAR-physics RS-MLLM published~\cite{sarchatchat2025}              & C \\
HSI reasoning        & Weak (all 18 models fail)      & Critical                & Very High   & HM-Bench: systematic failure~\cite{hmbench2026}                      & B \\
Temporal reasoning   & Emerging (Earth-OV temporal; VLRS-Bench) & Critical                & Very High   & Earth-OneVision~\cite{earthonevision2026}; VLRS-Bench~\cite{vlrsbench2026}     & C \\
Uncertainty output   & None (no calibrated output)    & Critical                & Very High   & Not addressed in any RS-MLLM                                         & D \\
Safety / alignment   & Emerging (1 paper)             & Critical                & Very High   & RemoteShield~\cite{remoteshield2026}; RLHF~\cite{rlhfrs2025}         & C \\
\botrule
\end{tabular*}
\footnotetext{``Operational requirement'' is assessed against the functional needs
of geospatial intelligence applications including disaster response, infrastructure
monitoring, and agricultural management---applications where incorrect spatial
output (wrong bounding box, missed object, incorrect land cover label) carries
real-world consequences.
``Gap'' ratings are qualitative assessments derived from the quantitative gaps
documented in Sects.~\ref{sec4}--\ref{sec7}; they represent the editors'
synthesis of the evidence, not formal meta-analytic effect sizes.}
\end{table*}

The matrix reveals a clear two-tier structure: optical VQA and captioning
are approaching operational capability (low gap), while OBB detection, SAR
reasoning, HSI reasoning, temporal reasoning, uncertainty quantification,
and safety are all in a ``critical gap'' state (very high gap). This
pattern directly motivates the eight research directions below.

\subsection{Critical Synthesis: Five Common Evaluation Fallacies}\label{sec9:fallacies}

The Operational Readiness Matrix reveals not only capability gaps but also
a systematic pattern of evaluation practices that overstate RS-MLLM capability.
We identify five evaluation fallacies that consistently inflate reported performance
and obscure genuine operational readiness.

\textbf{F1: Overall VQA score $\neq$ spatial reasoning.}
RSVQA-LR accuracy above 90\,\% is reported by multiple models, but RSVQA-LR
tests existence, counting, and comparison queries over low-resolution Sentinel-2
imagery. High RSVQA-LR scores do not imply the ability to locate objects, reason
about spatial relationships at sub-pixel resolution, or produce grounded bounding
boxes. Papers citing RSVQA-LR as evidence of RS understanding conflate
classification accuracy with spatial intelligence.

\textbf{F2: Single-dataset performance $\neq$ generalisation.}
The majority of RS-MLLM evaluations report results on one or two benchmarks
(typically RSVQA-LR/HR and RSICD), trained on instruction data from the same
distribution. Cross-dataset evaluation---testing a model trained on GeoChat-Instruct
on VRSBench, for instance---is rare. Without it, reported scores measure
in-distribution fitting, not generalisation.

\textbf{F3: Optical-only performance $\neq$ multimodal RS intelligence.}
A model achieving state-of-the-art on RSICD captioning has demonstrated optical
RS capability only. The 84.6\,\% optical dominance in the corpus means that most
published performance numbers reflect optical-only capability; they say nothing
about SAR interpretation, HSI spectral reasoning, or cross-sensor generalisation,
all of which are required for operational RS.

\textbf{F4: Textual answer accuracy $\neq$ spatial grounding.}
VQA tasks accept free-text answers matched against reference strings.
A model answering ``yes, there is a ship'' correctly scores the same as one that
could also localize the ship to sub-10-pixel precision. The VQA metric collapses
the distinction between perceptual identification and spatial grounding---a
distinction that is operationally critical for disaster response and maritime
surveillance.

\textbf{F5: Model scale $\neq$ geospatial intelligence.}
The shift from 7B to 13B or larger LLM backbones does not close the OBB gap,
the SAR physics gap, or the HSI spectral gap. These gaps are architectural in
origin---they require encoder redesign and physics-informed training, not parameter
scaling. Papers claiming geospatial intelligence improvements from backbone scaling
alone should be scrutinised carefully against the Operational Readiness Matrix
(Table~\ref{tab:orm}).

These five fallacies are not attributable to intentional misrepresentation but
to the natural incentive structure of benchmark optimisation. Addressing them
requires the field to adopt the richer evaluation protocol implied by the
Operational Readiness Matrix: cross-dataset evaluation, multimodal tasks,
spatial grounding metrics, and physics-informed probing.

\subsection{Research Directions}\label{sec9:directions}

Each direction is connected to the gap it addresses, the technical
approach it proposes, and---where applicable---to relevant methodological foundations
in the literature. Figure~\ref{fig:future_roadmap} presents
a roadmap organising the eight directions by implementation horizon and
estimated impact.

\begin{figure}[h]
\centering
\begin{tikzpicture}[font=\small, >=Stealth,
    box/.style={draw, rounded corners=3pt, fill=#1, minimum width=4.6cm,
                minimum height=0.72cm, align=center, font=\scriptsize},
    arr/.style={->, thick, gray!50}]

  %% Y-axis label
  \node[rotate=90, gray!65, font=\footnotesize] at (-0.7,2.5)
    {Estimated Impact};
  \draw[->,thick,gray!50] (0,0) -- (0,5.4);
  \foreach \y/\lbl in {0/Low, 2.7/Medium, 5.4/High}{
    \node[left,font=\tiny,gray!55] at (-0.05,\y) {\lbl};
  }
  %% X-axis
  \draw[->,thick,gray!50] (0,0) -- (11.2,0);
  \node[below,font=\footnotesize,gray!65] at (5.5,-0.5)
    {Implementation Horizon};
  \foreach \x/\lbl in {1.5/Near (1--2\,yr), 5.5/Mid (2--4\,yr), 9.5/Long (4+\,yr)}{
    \node[below,font=\scriptsize] at (\x,-0.2) {\lbl};
    \draw[gray!30,dashed] (\x,0) -- (\x,5.2);
  }

  %% ── Near-term (high impact) ─────────────────────────────
  \node[box=teal!18]   at (1.5,5.0) {D1: RS-Specific OBB\\MLLM Training};
  \node[box=teal!18]   at (1.5,4.2) {D2: Type~III Hallucination\\Benchmark};
  \node[box=orange!18] at (1.5,3.2) {D3: Physics-Informed\\SAR Encoder};

  %% ── Mid-term ─────────────────────────────────────────────
  \node[box=teal!22]   at (5.5,5.1) {D4: Unified Multi-Sensor\\RS-MLLM};
  \node[box=orange!22] at (5.5,4.1) {D5: Spectral-Physics\\HSI Encoder};
  \node[box=purple!18] at (5.5,3.0) {D6: Spatiotemporal\\Reasoning};

  %% ── Long-term ────────────────────────────────────────────
  \node[box=red!14]    at (9.5,5.0) {D7: RS Geospatial\\Foundation Model};
  \node[box=gray!20]   at (9.5,3.5) {D8: Operational RS-MLLM\\Safety Framework};

  %% Arrows connecting related directions
  \draw[arr] (3.8,5.0) -- (3.2,5.0);  % D1 feeds D4
  \draw[arr] (3.8,3.2) -- (3.2,3.2);  % D3 feeds D4
  \draw[arr] (7.8,5.1) -- (7.2,5.1);  % D4 feeds D7
  \draw[arr] (7.8,4.1) -- (7.2,4.1);  % D5 feeds D7

\end{tikzpicture}
\caption{Future research roadmap for RS-MLLMs, organised by implementation
horizon (near: 1--2 years; mid: 2--4 years; long: 4+ years) and estimated
impact on operational geospatial intelligence. Near-term directions (D1--D3)
address structural gaps identifiable within the current architectural paradigm;
long-term directions (D7--D8) require new paradigms and research
infrastructure.}\label{fig:future_roadmap}
\end{figure}

%% ----------------------------------------------------------
\subsection{D1: RS-Specific OBB-Capable MLLM Training}\label{sec9:d1}
%% ----------------------------------------------------------

\textbf{Gap addressed.} As established in Sects.~\ref{sec3:spatial}
and~\ref{sec4:obb}, four Corpus~B models (GeoChat, EarthGPT, GeoGround, EarthGPT-X)
support OBB output, but all are confined to optical DOTA~\cite{xia2018dota}
imagery; cross-sensor OBB reasoning is essentially unaddressed.
No RS-MLLM achieves OBB detection competitive with specialised detectors
such as Oriented R-CNN~\cite{xie2021orientedrcnn}.

\textbf{Research direction.} The fundamental challenge is extending MLLM
output heads from text-only generation to structured angular coordinate
prediction. Two complementary approaches are tractable in the near term.
First, a \emph{dedicated OBB output head} can be appended after the LLM
decoder, trained with an $\ell_1$ regression loss on the five OBB
parameters $(x, y, w, h, \theta)$ while the LLM backbone is kept
frozen\,+\,LoRA. This requires a training dataset with OBB annotations
in conversational format---currently sparse. Second, a
\emph{text-formatted OBB} approach encodes the five OBB parameters as
special tokens in the LLM output vocabulary, extending GeoChat's existing
design to multi-class, multi-sensor OBB settings.

\textbf{Methodological foundation.} Several candidate approaches exist for
constructing OBB-rich instruction datasets. Annotation efficiency methods
for long-tail RS object categories---including class-rebalancing
strategies~\cite{ding2021object} and semi-supervised pseudo-label
propagation~\cite{sm2lf2026}---provide a basis for scaling OBB annotation
to tail categories. Physics-informed augmentation~\cite{dsc_aug2025} and
optimisation-driven feature learning approaches~\cite{odfe2025} further
address the class imbalance and annotation scarcity problems identified
in Sect.~\ref{sec2:longtail}. Taken together, these methods provide
a concrete starting point for the OBB instruction data pipeline that
this direction requires.

%% ----------------------------------------------------------
\subsection{D2: A Comprehensive Type~III Hallucination Benchmark}\label{sec9:d2}
%% ----------------------------------------------------------

\textbf{Gap addressed.} Section~\ref{sec7:bench} established that no
evaluation benchmark provides comprehensive Type~III (Spectral Hallucination)
evaluation. HM-Bench~\cite{hmbench2026} addresses HSI tasks but does not
provide per-query hallucination labels or SAR-specific evaluation.

\textbf{Research direction.} A dedicated Type~III benchmark requires three
components: (i)~co-registered optical, SAR, and HSI image triplets of the
same geographic scenes, providing ground-truth sensor-physics information
independent of colour appearance; (ii)~per-modality query-response pairs
annotated with a hallucination label indicating whether the model's
response was driven by spectral physics or by RGB visual priors; and
(iii)~at least 5,000--10,000 triplets at diverse geographic locations and
land-cover types to enable statistically meaningful cross-modality analysis.
The DFC 2018 dataset (Sentinel-1 SAR, Sentinel-2 multispectral, and ISPRS
Vaihingen DSM)~\cite{yokoya2018hyperspectral} provides a starting point
for co-registered multi-sensor annotation.

\textbf{Estimated effort.} Constructing a 10,000-triplet benchmark with
expert annotation is a tractable 12--18 month effort for a well-staffed
research group. The resulting benchmark would directly enable evaluation of
all future SAR-aware and HSI-aware RS-MLLM architectures.

%% ----------------------------------------------------------
\subsection{D3: Physics-Informed SAR Encoder for RS-MLLMs}\label{sec9:d3}
%% ----------------------------------------------------------

\textbf{Gap addressed.} Section~\ref{sec5:sarbarriers} established that all
current RS-MLLM encoders are trained on RGB imagery and have no mechanism
for SAR backscatter physics interpretation. This causes Type~III spectral
hallucination on SAR inputs.

\textbf{Research direction.} A physics-informed SAR encoder should be
pre-trained with objectives derived from SAR scattering physics rather than
contrastive image-text alignment. Concretely, three complementary objectives
are proposed. First, a \emph{backscatter prediction objective} trains the
encoder to predict the normalised radar cross-section ($\sigma^0$) value
from SAR amplitude, providing a physical grounding absent from CLIP-style
contrastive training. Second, a \emph{scattering-mechanism classification
objective} trains the encoder to distinguish specular, double-bounce, and
volume scattering regimes within SAR patches, equipping the model with
the interpretive vocabulary needed for SAR scene description. Third, a
\emph{cross-sensor consistency objective} trains the encoder to produce
representations consistent between co-registered optical and SAR images
of the same scene, directly addressing the cross-sensor alignment problem.

\textbf{Methodological foundation.} Several prior works demonstrate the
feasibility of physics-informed SAR--optical alignment. SAR domain adaptation
approaches~\cite{schmitt2019fusion} and multi-modal RS image understanding
frameworks~\cite{zhang2024phi} establish that cross-sensor alignment improves
when domain-specific constraints are incorporated. Rayleigh-constrained
augmentation~\cite{rayleigh2026} shows that incorporating SAR backscatter
physics into the data-level fusion pipeline substantially improves cross-sensor
feature alignment; extending this principle from the data pipeline to the
encoder training objective is a concrete and tractable next step grounded
in established physical optics methodology.

%% ----------------------------------------------------------
\subsection{D4: Unified Multi-Sensor RS-MLLM}\label{sec9:d4}
%% ----------------------------------------------------------

\textbf{Gap addressed.} Section~\ref{sec5:fusion} established that the
cross-modal fusion matrix has five of nine cells unoccupied, with all
HSI-involving cells empty and SAR-optical fusion limited to a single
model without physics grounding.

\textbf{Research direction.} A unified multi-sensor RS-MLLM requires
three architectural innovations beyond the current state of the art.
First, a \emph{modality-specific encoder bank}: separate encoder branches
for optical (CLIP-based), SAR (physics-informed, as in D3), and HSI
(spectral-physics-aware, as in D5), each pre-trained with modality-appropriate
objectives. Second, a \emph{modality-aware connector}: an extension of
the ACSA connector~\cite{earthonevision2026} that explicitly conditions
token compression on the modality type, preventing optical-optimised
compression strategies from discarding spectral or backscatter features.
Third, a \emph{modality prompt token}: a learnable token prepended to
the LLM input that signals the active sensor modality, analogous to
language tokens in multilingual LLMs~\cite{touvron2023llama}, enabling
the model to switch its interpretation context without architectural
modification.

Earth-OneVision~\cite{earthonevision2026} is the closest existing
system to this direction, processing six sensor modalities within a
single model. However, it uses a unified optical-optimised encoder for
all modalities without physics-informed specialisation, and does not
include HSI. A true unified multi-sensor RS-MLLM would close both gaps.

%% ----------------------------------------------------------
\subsection{D5: Spectral-Physics-Aware HSI Encoder}\label{sec9:d5}
%% ----------------------------------------------------------

\textbf{Gap addressed.} Section~\ref{sec5:hsi} established that no
published RS-MLLM encoder handles HSI input correctly, and
HM-Bench~\cite{hmbench2026} confirms systematic Type~III spectral
hallucination across 18 tested models.

\textbf{Research direction.} An HSI-compatible encoder must address the
dimensional incompatibility problem (HSI data cubes have $B\!\in\![200,400]$
spectral bands versus the 3-channel RGB assumption of CLIP) without
discarding the spectral information that makes HSI valuable.

Two encoder designs are tractable. First, a \emph{spectral tokeniser}:
a 1D convolutional network that maps the full spectral signature of each
spatial pixel to a fixed-length feature vector, applied independently
at each spatial position before spatial ViT processing. This preserves
spectral structure across the full $B$-band range while producing a
representation compatible with standard MLLM connectors. Second, a
\emph{band-grouping ViT}: the $B$ spectral bands are partitioned into
physically meaningful groups (visible, NIR, SWIR, etc.) and processed
by separate ViT branches whose outputs are fused, aligning the encoder's
inductive bias with the spectral structure of the data. Both designs
should be pre-trained with spectral unmixing objectives that predict
endmember abundances, equipping the encoder with the spectral
discrimination capabilities documented as absent in
HM-Bench~\cite{hmbench2026}.

%% ----------------------------------------------------------
\subsection{D6: Spatiotemporal Reasoning in RS-MLLMs}\label{sec9:d6}
%% ----------------------------------------------------------

\textbf{Gap addressed.} Section~\ref{sec4:reason} established that
Level-4 spatiotemporal reasoning---tracking object state changes,
inferring causal relationships across temporal acquisitions, answering
multi-hop questions---is effectively absent from published RS-MLLMs.

\textbf{Research direction.} Spatiotemporal reasoning requires three
capabilities beyond current RS-MLLM design: (i)~\emph{temporal token
fusion}, in which visual tokens from multiple temporal acquisitions
of the same area are jointly encoded with explicit temporal position
embeddings, analogous to temporal encoding in video transformers;
(ii)~\emph{change-aware attention}, in which cross-attention between
temporal tokens is guided by a pre-computed change mask, focusing the
model's attention budget on the regions that differ across time; and
(iii)~\emph{causal reasoning training data}, in which instruction pairs
explicitly require multi-step inference connecting observed changes to
physical or anthropogenic causes.

Temporal consistency is an active area in RS: TEOChat~\cite{teochat2025}
and Change-Agent~\cite{changeagent2025} represent the most advanced existing
work in conversational temporal RS reasoning. Semi-supervised temporal
consistency approaches---including pseudo-label agreement across temporal
frames~\cite{sm2lf2026} and video-MLLM architectures~\cite{llava_onevision2024}---provide
methodological foundations for extending these to causal multi-temporal
reasoning. Extending them to causal reasoning is the most tractable near-term step.
Two post-cutoff works corroborate this gap: LongEarth-R1~\cite{longearth_r1_2026}
introduces LongEarth-Bench ($\sim$120k QA, 12 long-sequence EO reasoning tasks),
and RSVideo~\cite{rsvideo_2026} introduces RSVideo-10K and RSVideo-Bench for
continuous RS video understanding---both confirming that long-horizon spatiotemporal
reasoning is the most active open frontier at the time of this survey's submission.

%% ----------------------------------------------------------
\subsection{D7: RS Geospatial Foundation Model}\label{sec9:d7}
%% ----------------------------------------------------------

\textbf{Gap addressed.} The current RS-MLLM landscape consists of
dozens of task-specific or domain-specific models, each optimised for
a narrow capability (VQA, captioning, OBB detection, or specific sensor
type). This fragmentation prevents the cross-task and cross-sensor
generalisation required for operational geospatial intelligence.

\textbf{Research direction.} A geospatial foundation model for RS
should be pre-trained on a diverse corpus of RS data across all sensor
modalities, spatial resolutions, geographic regions, and temporal
scales, with a training objective that jointly optimises for multiple
downstream tasks. Prithvi-100M~\cite{prithvi2023} and its multi-sensor
successor Prithvi-EO~\cite{geofm2024} represent early steps in this direction for
multitemporal optical RS; extending the foundation model paradigm to include SAR
and HSI modalities with physics-informed objectives is the critical next step.

The scale required for a true RS geospatial foundation model---estimated
at 100M--1B parameters with training data in the 100M+ sample
range---is beyond current academic compute budgets but is achievable
through national-scale research infrastructure or industrial partnerships.
The InternVL~\cite{internvl2024} and GeoLLM~\cite{geollm2025} architectures demonstrate that scaling
vision-language foundation models with careful multi-task training
produces emergent cross-task generalisation; an analogous approach
applied to RS-specific data with physics-aware objectives represents
the long-term target for the field.

%% ----------------------------------------------------------
\subsection{D8: Operational RS-MLLM Safety Framework}\label{sec9:d8}
%% ----------------------------------------------------------

\textbf{Gap addressed.} Section~\ref{sec7:danger} established that
RS-MLLM hallucinations carry qualitatively higher stakes than their
natural-image counterparts, with direct consequences for disaster
response, infrastructure monitoring, and agricultural policy. The
field currently lacks a systematic safety framework for operational
RS-MLLM deployment.

\textbf{Research direction.} An operational RS-MLLM safety framework
should include four components. First, \emph{calibrated confidence
estimation}: RS-MLLMs should output calibrated uncertainty scores
alongside their responses, enabling downstream systems to flag low-confidence
outputs for human review. Second, \emph{domain-specific hallucination
detection}: automated detection of Type~I, Type~II, and Type~III
hallucinations during inference, building on the benchmark infrastructure
of RSHallu~\cite{rshallu2026} and the Type~III benchmark proposed in D2.
Third, \emph{cross-sensor consistency checking}: for applications with
access to multiple sensor modalities, inconsistencies between RS-MLLM
responses to different sensor views of the same scene should trigger
automatic flagging. Fourth, \emph{operational evaluation protocols}:
standardised benchmarks that measure RS-MLLM performance specifically
in operational contexts (disaster response, infrastructure inspection,
precision agriculture), complementing the general VQA and captioning
benchmarks that dominate current evaluation.

\subsection*{Summary}

Table~\ref{tab:future_summary} summarises the eight future research
directions, cross-referenced to the gaps they address and the
prior-work foundations that make each direction tractable.

\begin{table}[h]
\caption{Summary of eight future research directions for RS-MLLMs.
``Gap'' column references the relevant section; ``Foundation'' lists
the methodological building block most directly applicable.
Cited prior work (ODFE~\cite{odfe2025}, SM$^2$-LF~\cite{sm2lf2026}, Rayleigh~\cite{rayleigh2026}) provides methodological
foundations that make these directions tractable; these are starting
points, not completed solutions.}\label{tab:future_summary}
\small\setlength\tabcolsep{4pt}
\begin{tabular*}{\textwidth}{@{\extracolsep\fill}lp{3.0cm}p{3.2cm}p{3.5cm}}
\toprule
\textbf{Dir.} & \textbf{Title} & \textbf{Gap (Sect.)} & \textbf{Foundation} \\
\midrule
D1 & OBB-Capable MLLM
   & Sec.~\ref{sec3:spatial}, Sec.~\ref{sec4:obb}
   & Ding et al.~\cite{ding2021object}; ODFE~\cite{odfe2025}; SM$^2$-LF~\cite{sm2lf2026} \\[3pt]
D2 & Type~III Hallucin.\ Bench
   & Sec.~\ref{sec7:bench}
   & HM-Bench~\cite{hmbench2026}; DFC 2018~\cite{yokoya2018hyperspectral} \\[3pt]
D3 & Physics SAR Encoder
   & Sec.~\ref{sec5:sarbarriers}, Sec.~\ref{sec7:type3}
   & Schmitt et al.~\cite{schmitt2019fusion}; Rayleigh~\cite{rayleigh2026} \\[3pt]
D4 & Unified Multi-Sensor MLLM
   & Sec.~\ref{sec5:fusion}
   & Earth-OneVision~\cite{earthonevision2026}; D3, D5 \\[3pt]
D5 & Spectral HSI Encoder
   & Sec.~\ref{sec5:hsibarriers}
   & HM-Bench~\cite{hmbench2026}; bioucas2013~\cite{bioucas2013hyperspectral} \\[3pt]
D6 & Spatiotemporal Reasoning
   & Sec.~\ref{sec4:reason}
   & TEOChat~\cite{teochat2025}; Change-Agent~\cite{changeagent2025}; SM$^2$-LF~\cite{sm2lf2026} \\[3pt]
D7 & RS Foundation Model
   & Sec.~\ref{sec3:master}
   & Prithvi-100M~\cite{prithvi2023}; InternVL~\cite{internvl2024} \\[3pt]
D8 & Safety Framework
   & Sec.~\ref{sec7:danger}
   & RSHallu~\cite{rshallu2026}; RemoteShield~\cite{remoteshield2026}, RLHF approaches~\cite{rlhfrs2025}, \\
\botrule
\end{tabular*}
\end{table}

The eight directions form two priority tiers. Directions D1--D3 are
near-term: they address structural gaps (OBB output, Type~III benchmark,
SAR encoder) that can be closed within the current MLLM architectural
paradigm with targeted engineering effort. Directions D4--D6 are mid-term:
they require new architectures and training objectives but build directly
on the near-term work. Directions D7--D8 are long-term: D7 requires a
scale of compute and data not currently available in academia; D8 requires
community consensus on operational evaluation protocols that does not yet
exist. Together, the eight directions chart a concrete path from current
RS-MLLMs---perception-limited systems that describe optical imagery
conversationally---to geospatially intelligent systems that reason
operationally across all sensor modalities, spatial scales, and temporal
acquisitions.

%% ──────────────────────────────────────────────────────────────
\subsection{Limitations of This Survey}\label{sec9:limitations}
%% ──────────────────────────────────────────────────────────────

We acknowledge nine limitations of this systematic review; transparency about
these limitations increases rather than decreases the trustworthiness of the synthesis.

\textbf{L1: Protocol pre-registration.} The review protocol was not
pre-registered in PROSPERO or an equivalent registry before search
execution. This is a common limitation of rapidly evolving technology
surveys, where the literature landscape changes faster than the standard
pre-registration cycle allows. We report our methods in full detail
(Sect.~\ref{sec2:prisma}, Tables~\ref{tab:inclusion_exclusion},
\ref{tab:search_strategy}, \ref{tab:quality}, \ref{tab:prisma_checklist})
to compensate for the absence of pre-registration.

\textbf{L2: Rapidly changing field.} The RS-MLLM field is evolving
rapidly, with multiple new RS-MLLM and benchmark releases appearing
throughout 2026. Any specific quantitative finding (model rankings,
benchmark scores) should be
interpreted as a snapshot of the August 2026 state of the field rather
than a durable ranking. Sensitivity analysis SA2 (Sect.~\ref{sec2:prisma})
confirms that all structural conclusions hold across the 2022--2025
window.

\textbf{L3: Preprint-heavy literature.} arXiv preprints account for
92 of the 210 total retrieved papers (44\,\%). Preprints have not completed peer
review, and reported results may not be independently replicated. We
applied the I2 reproducibility gate (public code or weights required)
to mitigate this, but benchmark scores from preprints should be
interpreted with greater caution than those from peer-reviewed papers.

\textbf{L4: Database coverage.} Searches covered 8 databases/proceedings.
Chinese-language conferences (e.g., CCVS, CSIG), niche RS proceedings,
and institutional technical reports are not covered, which may introduce
coverage bias toward English-language and Western-venue publications.

\textbf{L5: Terminology inconsistency.} The RS-MLLM field uses
inconsistent terminology: ``MLLM'', ``VLM'', ``vision-language model'',
``multimodal LLM'', and ``foundation model'' are used interchangeably
across papers. Our Boolean search covered all major variants; however,
novel terminological variants introduced after August 2026 may not be
captured.

\textbf{L6: Heterogeneous benchmarks.} Performance comparisons across
papers are limited by benchmark heterogeneity: different papers use
different test splits, prompt formats, and evaluation scripts for the
same benchmark. Table~\ref{tab:cvmllm_comparison} (Sect.~\ref{sec6:cvmllm})
distinguishes same-protocol from cross-paper comparisons and marks the
latter as not directly comparable.

\textbf{L7: Incomplete reporting.} 44\,\% of included papers do not
explicitly acknowledge limitations (quality criterion Q7, Table~\ref{tab:quality}).
Architecture details were recorded as ``NR'' (not reported) when absent
from both the paper and its associated code repository; these gaps
reduce the completeness of the quantitative synthesis.

\textbf{L8: Publication and preprint bias.} Papers with positive or
strong results are more likely to be submitted and accepted. Performance
figures reported in the synthesis may therefore be optimistic relative
to the true distribution of RS-MLLM capability. We note that 15\,\%
of included papers are classified as Low quality (Table~\ref{tab:quality}),
providing some evidence of the lower tail of the distribution.

\textbf{L9: Evolving model versions.} Several included models released
updated versions after the search cutoff (e.g., Earth-OneVision, GeoChat).
We report results for the version available at the time of search;
later versions may differ substantially in performance.

\textbf{L10: Cross-paper metric incomparability.} Even where the same
benchmark is used, different papers employ different test splits, prompt
formats, evaluation scripts, and answer normalisation procedures. As a
result, numerical scores from different papers on the same benchmark
(e.g., RSVQA-LR accuracy) are frequently not directly comparable.
Table~\ref{tab:cvmllm_comparison} (Sect.~\ref{sec6:cvmllm}) explicitly
distinguishes same-protocol from cross-paper comparisons and marks the
latter as not directly comparable; readers should exercise the same
caution when interpreting all performance figures in this survey.

\textbf{L11: Cutoff-date limitation.} The primary search cutoff is
14 August 2026; a targeted update search was conducted on 24 August 2026
(Sect.~\ref{sec2:update_search}), identifying three post-cutoff works
(LongEarth-R1~\cite{longearth_r1_2026}, RSVideo~\cite{rsvideo_2026},
and the IGARSS 2026 RS-VLM session papers) consistent with but not
altering any survey finding. Papers submitted after 24 August 2026 are
not included. Given the RS-MLLM field's rapid growth rate (42 papers in
2025 alone), further relevant works will have appeared between the update
search and final publication. The structured quantitative findings
(Sect.~\ref{sec8}) should be interpreted as a 24 August 2026 snapshot;
the five-dimensional design-space taxonomy and hallucination framework
(Sects.~\ref{sec3},~\ref{sec7}) are structural and expected to remain
valid across the next generation of RS-MLLM development.
